\clearpage
\phantomsection
\addcontentsline{toc}{chapter}{TÀI LIỆU THAM KHẢO}
\begin{thebibliography}{99}

\bibitem{thanh2024}
Dương Hữu Thành, \textit{Công nghệ phần mềm}, NXB Thông tin và Truyền thông, 2024.

\bibitem{duong2026}
H. T. Duong, T. P. Bui, D. H. Nguyen, ``QLoRA and RAG Application in Building
Admissions Support Chatbot: Performance Evaluation and Knowledge Update
Ability,'' in N. N. Dao, V. Truong Hoang, F. Dornaika (eds.),
\textit{Explainable Intelligence in Digital Twins}, EIDT 2025, LNEE vol.~1531,
Springer, Singapore, 2026.

\bibitem{lewis2020rag}
P. Lewis, E. Perez, A. Piktus \textit{et al.}, ``Retrieval-Augmented Generation
for Knowledge-Intensive NLP Tasks,'' \textit{NeurIPS}, vol.~33, pp.~9459--9474,
2020.

\bibitem{karpukhin2020dpr}
V. Karpukhin, B. Oğuz, S. Min \textit{et al.}, ``Dense Passage Retrieval for
Open-Domain Question Answering,'' \textit{Proc. EMNLP}, pp.~6769--6781, 2020.

\bibitem{robertson2009bm25}
S. Robertson, H. Zaragoza, ``The Probabilistic Relevance Framework: BM25 and
Beyond,'' \textit{Foundations and Trends in Information Retrieval}, vol.~3,
no.~4, pp.~333--389, 2009.

\bibitem{cormack2009rrf}
G. V. Cormack, C. L. A. Clarke, S. Buettcher, ``Reciprocal Rank Fusion
Outperforms Condorcet and Individual Rank Learning Methods,'' \textit{Proc.
SIGIR}, pp.~758--759, 2009.

\bibitem{gao2023hyde}
L. Gao, X. Ma, J. Lin, J. Callan, ``Precise Zero-Shot Dense Retrieval without
Relevance Labels,'' \textit{Proc. ACL}, pp.~1762--1777, 2023.

\bibitem{quochoi2017}
Quốc hội, \textit{Luật Hỗ trợ doanh nghiệp nhỏ và vừa}, số 04/2017/QH14.

\bibitem{bll2019}
Quốc hội, \textit{Bộ luật Lao động}, số 45/2019/QH14.

\bibitem{qlthue2019}
Quốc hội, \textit{Luật Quản lý thuế}, số 38/2019/QH14.

\bibitem{postgresqlfts}
PostgreSQL Global Development Group, ``Full Text Search,'' tài liệu trực tuyến,
\url{https://www.postgresql.org/docs/current/textsearch.html}.

\bibitem{rfc7519}
M. Jones, J. Bradley, N. Sakimura, ``JSON Web Token (JWT),'' RFC~7519, 2015.

\bibitem{fastapi}
S. Ramírez, \textit{FastAPI documentation}, tài liệu trực tuyến,
\url{https://fastapi.tiangolo.com}.

\bibitem{react}
Meta, \textit{React documentation}, tài liệu trực tuyến,
\url{https://react.dev}.

\bibitem{docker}
Docker Inc., \textit{Docker documentation}, tài liệu trực tuyến,
\url{https://docs.docker.com}.

\bibitem{chromadb}
Chroma, \textit{Chroma documentation}, tài liệu trực tuyến,
\url{https://docs.trychroma.com}.

\bibitem{rfc9110}
R. Fielding, M. Nottingham, J. Reschke, ``HTTP Semantics,'' RFC~9110, 2022.

\bibitem{rfc8259}
T. Bray, ``The JavaScript Object Notation (JSON) Data Interchange Format,''
RFC~8259, 2017.

\bibitem{vaswani2017}
A. Vaswani, N. Shazeer, N. Parmar \textit{et al.}, ``Attention Is All You Need,''
\textit{NeurIPS}, 2017.

\bibitem{argon2}
A. Biryukov, D. Dinu, D. Khovratovich, ``Argon2: the memory-hard function for
password hashing and other applications,'' winner of the Password Hashing
Competition, 2015.

\bibitem{pdt590}
pdt590, ``vietnamese-legal-documents,'' Hugging Face Datasets, 2024.
\url{https://huggingface.co/datasets/pdt590/vietnamese-legal-documents}.

\end{thebibliography}

\clearpage
\phantomsection
\addcontentsline{toc}{chapter}{PHỤ LỤC}
\chapter*{PHỤ LỤC}

\section*{Phụ lục A. Chạy hệ thống trên máy}

Sao \texttt{.env.example} thành \texttt{.env} \textbf{ở thư mục gốc dự án},
không phải trong \texttt{backend/}. Compose không đọc \texttt{backend/.env}.
Điền \texttt{LLM\_API\_KEY} (trang AI Studio của Google) và
\texttt{JWT\_SECRET\_KEY} (ví dụ \texttt{openssl rand -hex 32}). Chạy:

\texttt{docker compose up -d --build}

Lần đầu nạp luật:

\texttt{docker compose exec backend python scripts/load\_postgres.py}
\texttt{--truncate}

Rồi restart backend. Vào \texttt{http://localhost:5173}. Tài khoản đăng ký đầu
tiên thành quản trị. Tra cứu dùng được ngay. Hỏi đáp cần đã nhúng vector: vào
\texttt{/admin/corpus} bấm reindex, hoặc đợi startup nếu khóa đã có từ đầu.

Thiếu khóa: web vẫn mở, hỏi đáp trả 503 chữ tiếng Việt, không phải stack trace.

\section*{Phụ lục B. Một số đường API}

\begin{longtable}{p{6.6cm}p{7.2cm}}
\toprule
\textbf{Đường} & \textbf{Việc} \\
\midrule
\endfirsthead
\toprule
\textbf{Đường} & \textbf{Việc} \\
\midrule
\endhead
\texttt{POST /api/v1/auth/register} & Tạo user, có thể kèm organization \\
\texttt{POST /api/v1/auth/login} & Cấp access + refresh \\
\texttt{POST /api/v1/auth/refresh} & Đổi access, kiểm tra \texttt{type} \\
\texttt{GET /api/v1/conversations} & Danh sách hội thoại của đúng user \\
\texttt{POST /api/v1/legal/chat/stream} & SSE hỏi đáp \\
\texttt{GET /api/v1/laws} & Catalog \\
\texttt{GET /api/v1/laws/search} & FTS \\
\texttt{GET /api/v1/laws/\{id\}/articles} & Cây Chương--Điều \\
\texttt{POST /api/v1/documents} & Tải file \\
\texttt{POST /api/v1/documents/\{id\}/review} & 202, chạy nền \\
\texttt{GET /api/v1/compliance/tasks} & Lịch \\
\texttt{POST /api/v1/admin/corpus/reindex} & Nhúng lại, chỉ admin \\
\bottomrule
\end{longtable}

\section*{Phụ lục C. Chạy hàng loạt câu hỏi}

Quản trị mở trang \texttt{/admin/batch-eval}, gửi một mảng JSON (mỗi phần tử có
câu hỏi) để pipeline chạy lần lượt, xem tiến độ và tải kết quả. Chế độ này bỏ
qua phân tích ý định (luôn truy hồi) và không gắn lịch sử hội thoại, nhằm tái
lập được trên cùng bộ câu. User thường không gọi được các đường này.

\section*{Phụ lục D. Ghi chú biên dịch báo cáo}

File \texttt{main.tex} phải chạy XeLaTeX (gói \texttt{fontspec}). Trong Cursor
hoặc VS Code, biên dịch bằng recipe xelatex (hoặc \texttt{latexmk -xelatex}).
Đầu file có chỉ thị TEX program trỏ sang xelatex. Ảnh minh họa nằm trong thư mục
\texttt{Baocao/images/}; không dùng tùy chọn \texttt{demo} của \texttt{graphicx}
(tùy chọn đó vẽ hình đen đặc).

\section*{Phụ lục E. Biến môi trường chính}

\begin{longtable}{p{5.2cm}p{8.6cm}}
\toprule
\textbf{Biến} & \textbf{Ý nghĩa} \\
\midrule
\endfirsthead
\toprule
\textbf{Biến} & \textbf{Ý nghĩa} \\
\midrule
\endhead
\texttt{LLM\_API\_KEY} & Khóa gọi Gemini (sinh câu) \\
\texttt{EMBEDDINGS\_API\_KEY} & Khóa nhúng; có thể trùng khóa LLM \\
\texttt{JWT\_SECRET\_KEY} & Khóa ký JWT, tối thiểu đủ dài ngẫu nhiên \\
\texttt{LEGAL\_DATABASE\_URL} & Chuỗi kết nối Postgres (Compose tự ghép) \\
\texttt{CORS\_ALLOW\_ORIGINS} & Danh sách origin, dùng lúc dev Vite \\
\texttt{LLM\_MODEL} & Mặc định \texttt{gemini-2.5-flash} \\
\texttt{EMBEDDINGS\_MODEL} & Mặc định \texttt{gemini-embedding-001} \\
\texttt{POSTGRES\_PORT} & Mặc định 23432 trên host \\
\bottomrule
\end{longtable}

\section*{Phụ lục F. Cấu trúc thư mục mã nguồn}

\begin{itemize}
    \item \path{backend/src/routers}: cổng HTTP.
    \item \path{backend/src/services/agents}: LangGraph.
    \item \path{backend/src/services/vector_store}: BM25, Chroma, hybrid.
    \item \path{backend/src/services/legal}: FTS, catalog.
    \item \path{backend/src/services/contracts}: trích chữ, soát xét.
    \item \path{backend/src/services/compliance}: seed và sinh lịch.
    \item \path{backend/src/models}: ORM.
    \item \path{backend/alembic}: migration.
    \item \path{frontend/src/pages}: trang.
    \item \path{frontend/src/components/chat}: hội thoại.
    \item \path{corpus/law_manifest.json}: danh mục số hiệu.
    \item \path{data/base_data.json}: Điều.
    \item \path{Baocao/}: nguồn báo cáo XeLaTeX.
\end{itemize}

\section*{Phụ lục G. Use case còn lại (rút gọn)}

UC-01: khách gửi email, mật khẩu, hồ sơ công ty; hệ thống tạo user, có thể gắn
admin nếu là người đầu; sinh lịch nếu có tổ chức. UC-02: login cấp hai token;
refresh đổi access. UC-04: đổi tên hội thoại, xóa mềm hoặc xóa hẳn tùy API; không
xem hội thoại người khác. UC-06: ô tìm trên \texttt{/laws}, ba tầng FTS. UC-08:
thành viên cùng tổ chức xem finding. UC-10: admin import JSON, reindex. UC-11:
admin khóa tài khoản hoặc đổi vai trò. UC-12: admin gửi bộ câu, nhận file kết quả.

\section*{Phụ lục H. Thuật ngữ gắn với cách đóng gói đồ án}

\textbf{Container} ở đây là một trong ba tiến trình Compose (Postgres, backend,
nginx). \textbf{Image} là bản đóng gói từng dịch vụ. \textbf{Volume}
\texttt{postgres\_data}, \texttt{chroma\_db}, \texttt{uploads} giữ dữ liệu khi
tắt cụm. \textbf{Proxy} là nginx nhận một cổng rồi chuyển \texttt{/api} vào
backend. \textbf{Healthcheck} là \texttt{pg\_isready} và \texttt{/health}, để
backend đợi CSDL sẵn sàng. \textbf{Migration} là bản Alembic đổi cột.
\textbf{Interceptor} là hàm Axios gắn token và gỡ JSON khi gửi file.
\textbf{Splat} là đoạn URL bắt phần còn lại, dùng cho số hiệu có dấu \texttt{/}.

\section*{Phụ lục I. Quy trình nạp corpus và lập chỉ mục}

Nạp chữ vào Postgres tách khỏi nhúng vector. \texttt{scripts/load\_postgres.py}
đọc JSON Điều, ghi \texttt{legal\_knowledge\_records} (tuỳ chọn
\texttt{--truncate}). Generated column tự tính số Điều và tsvector. Tra cứu dùng
ngay sau bước này, không cần khóa Gemini.

Nhúng: lúc khởi động hoặc khi quản trị bấm reindex, hệ thống đọc từng Điều, gọi
embedding theo lô 64 bản ghi, ghi vào Chroma, rồi ghi manifest (mô hình, số
chiều, số vector). Khác mô hình hoặc khác chiều so với
manifest: bắt buộc nhúng lại. Thiếu khóa: bỏ qua nhúng, \texttt{/health} vẫn 200.

BM25: tách từ, cache đĩa theo tokenizer. Không nhầm cache BM25 với volume
Chroma.

Thứ tự an toàn khi máy trắng: Compose up, đợi health, load\_postgres, restart
backend, (nếu có khóa) reindex trên admin. Đảo thứ tự (reindex trước khi có
hàng Postgres) cho chỉ mục rỗng hoặc lệch id.

\section*{Phụ lục J. Ràng buộc đạo đức và trách nhiệm sử dụng}

Hệ thống hỗ trợ thông tin trên kho đóng. Không thay thế luật sư, đại lý thuế,
hoặc cơ quan nhà nước. Người dùng phải mở Điều gốc trước khi hành động có hậu
quả pháp lý. Nhật ký hội thoại chứa câu hỏi doanh nghiệp --- triển khai thật
cần chính sách lưu trữ và quyền xóa (đồ án mới có xóa hội thoại theo chủ sở
hữu). Không thu thập dữ liệu cá nhân ngoài tài khoản thử nghiệm.


